\section{Outlook: Extension to Multi-Qubit Systems}
\label{sec:block-encoding}

The single-qubit theory of Sections~\ref{sec:encoding}
and~\ref{sec:main} transforms a scalar signal $r$ into an
unbounded rational function via stereographic encoding, QSP,
and Pauli-measurement decoding.  A natural question is whether
this machinery extends to multi-qubit Hamiltonians, where the
goal is to apply a function to each eigenvalue simultaneously.

In the standard QSVT framework~\cite{GSLW2019}, a Hamiltonian
$H$ with $\|H\| = \alpha$ is embedded into a larger unitary via
an $(\alpha, m, \epsilon)$-block encoding, which stores
$H/\alpha$ in the top-left block.  The normalization factor
$\alpha \geq \|H\|$ inflates the circuit depth and forces the
QSP polynomial to approximate $f(\alpha\,\cdot\,)$ on $[-1,1]$.
The stereographic alternative removes this normalization by
encoding each eigenvalue directly onto the Bloch sphere of a
single ancilla qubit.

\begin{definition}[Stereographic Block Encoding]
	\label{def:stereo-block}
	Let $H$ be a Hermitian operator with spectral decomposition
	$H = \sum_j \lambda_j \ket{j}\bra{j}$, where
	$\lambda_j \geq 0$.  A stereographic block encoding of~$H$
	is a unitary $U$ acting on $n + 1$ qubits such that
	\begin{equation}
		U = \sum_j U_{\lambda_j} \otimes \ket{j}\bra{j},
		\label{eq:stereo-block}
	\end{equation}
	where $U_{\lambda_j}$ is the single-qubit encoding
	unitary~\eqref{eq:encoding-unitary} with $r = \lambda_j$
	and $\theta = 0$.
\end{definition}

In each eigenspace $\ket{j}$, the ancilla is rotated to the
encoding state
$\ket{\lambda_j} = (\lambda_j\ket{0} +
	\ket{1})/\sqrt{1+\lambda_j^2}$
via $R_y(2\gamma_j)$ with $\gamma_j = \arctan(1/\lambda_j)$.
The encoding is exact, uses a single ancilla qubit, and imposes
no constraint on $\|H\|$.  Once in place, the QSP protocol from
Section~\ref{subsec:qsp-phases} applies independently in each
eigenspace: a degree-$d$ sequence with phases $\bm{\phi}$
produces
$P(\tilde{\lambda}_j)\ket{0} +
	Q(\tilde{\lambda}_j)\ket{1}$
on the ancilla, and Pauli-measurement decoding yields the
rational eigenvalue transformation
\begin{equation}
	f(\lambda_j) = \frac{P(\tilde{\lambda}_j)}{Q(\tilde{\lambda}_j)},
	\label{eq:eigenvalue-transform}
\end{equation}
with circuit depth $\mathcal{O}(d)$ independent of $\|H\|$.
The positivity assumption $\lambda_j \geq 0$ can be relaxed by
a spectral shift $H \to H + |\lambda_{\min}|I$.

Constructing~\eqref{eq:stereo-block}, however, requires the
ability to apply $U_{\lambda_j}$ conditioned on the eigenspace
$\ket{j}$---it requires access to the eigenbasis.  This is a
strong assumption that standard block encoding does not make: in
the QSVT framework, $H/\alpha$ is embedded into a unitary
without knowledge of the eigenstates, and the QSP polynomial
acts simultaneously in all eigenspaces precisely because the
block encoding does not distinguish them.  The eigenbasis
requirement restricts stereographic block encoding to
\emph{structured Hamiltonians} whose eigenbasis is known or
efficiently computable---diagonal operators, integrable models,
or operators whose spectral decomposition is available as a
quantum circuit.  For generic Hamiltonians, the eigenbasis is
the object one wishes to compute, and the construction becomes
circular.  This limitation is not inherent to stereographic
encoding itself---the single-qubit theory holds for any scalar
signal---but rather to the specific block encoding strategy
of~\eqref{eq:stereo-block}.  Whether alternative strategies can
relax this requirement is an open question that we defer to a
companion paper.

To see concretely how the construction works despite this
restriction, consider the isotropic Heisenberg model with a
longitudinal field on two qubits:
\begin{equation}
	H = J(X \otimes X + Y \otimes Y + Z \otimes Z) + h(Z \otimes I),
	\label{eq:heisenberg}
\end{equation}
where $J$ is the exchange coupling and $h$ the field strength.
This is a non-trivial example requiring a genuine change of
basis: $\ket{00}$ and $\ket{11}$ are eigenstates with
eigenvalues $J + h$ and $J - h$, but the
$\{\ket{01}, \ket{10}\}$ subspace has eigenvalues
$\lambda_{2,3} = -J \pm \sqrt{h^2 + 4J^2}$
and eigenstates that reduce to Bell states when $h = 0$.  The
change-of-basis unitary $V$ acts non-trivially only on this
two-dimensional subspace as a Givens rotation, decomposable into
$\mathcal{O}(1)$ elementary gates.  The full stereographic block
encoding on 3 qubits (2 system $+$ 1 ancilla) is
\begin{equation}
	U = (V^\dagger \otimes I_a) \cdot U_{\mathrm{diag}} \cdot (V \otimes I_a),
	\label{eq:heisenberg-stereo}
\end{equation}
where $U_{\mathrm{diag}}$ is a uniformly controlled
$R_y(2\gamma_j)$ on the ancilla, conditioned on the
computational-basis state of the diagonalized system register:
\begin{equation}
	\begin{quantikz}[column sep=0.4cm, row sep=0.5cm]
		\lstick{$\ket{0}_a$} & \qw & \gate[style={inner sep=4pt}]{R_y(2\gamma_j)} & \qw & \qw \\
		\lstick{$\ket{\psi_1}$} & \gate[2]{V} & \ctrl{-1}\vqw{1} & \gate[2]{V^\dagger} & \qw \\
		\lstick{$\ket{\psi_0}$} & \qw & \ctrl{-1} & \qw & \qw
	\end{quantikz}
	\label{eq:heisenberg-circuit}
\end{equation}
The total gate count is constant and independent of $\|H\|$.
For $J = 1$, $h = 0.5$, the eigenvalues are
$\lambda \approx \{-3.06, 0.50, 1.06, 1.50\}$; after shifting
by $|\lambda_{\min}|$, the stereographic encoding compresses
even the largest shifted eigenvalue onto the Bloch sphere
without truncation.  This example is verified numerically in
Section~\ref{sec:cost-analysis}.

The Heisenberg model illustrates the general pattern: for any
Hamiltonian with a known diagonalization $H = V^\dagger D V$,
the stereographic block encoding reduces to a uniformly
controlled rotation sandwiched between $V$ and $V^\dagger$.
The difficulty is entirely in the diagonalization step---once
$V$ is available as a quantum circuit, the encoding is
straightforward.  A companion paper will develop this
construction systematically, analyze the decoding protocol when
the system register is in a superposition of eigenspaces, and
investigate the gate complexity for broader classes of
Hamiltonians.
